\subsection{重み付き平均}
% 最多区間内で最も近い2点以上を取得し，その平均値を基準とする．
% 全滞在者の時刻に平均値からの距離の逆数を重みとして付与し，積の和を求める．
% 重みの和で割ったものを帰宅推奨時刻としている．
% 最多区間内に含まれる帰宅予測時刻のうち，互いに近接する2点以上を抽出し，それらの平均値を基準時刻として算出する．
% 次に，全滞在者の帰宅予測時刻に対し，基準時刻からの時間的距離の逆数を重みとして付与し，重み付き平均を求めて帰宅推奨時刻を算出する．
% この方法により，基準時刻に近い利用者の影響を強く反映しつつ，離れた時刻の利用者の影響を緩やかに低減できる．
最多区間内に含まれる帰宅予測時刻を基に，複数人が同時に帰宅しやすい中心的な時刻を求め，それを軸として帰宅推奨時刻を算出する．
本手法では，最多区間内で時間的に近接したユーザ群を基準としつつ，滞在者の帰宅予測時刻を考慮し，一部の外れ値に左右されにくく同時帰宅が成立しやすい時刻の導出を目的とする．
以下では，まず基準時刻の定義方法を示し，次にその基準時刻を用いた帰宅推奨時刻の算出方法を数式で説明する．


% 式

最多区間に含まれる帰宅予測時刻の列の集合を
\[
  T' = \{ t'_1, t'_2, \dots, t'_M \} \quad (M \geq 2)
\]
とする．
このうち，最多区間内において時間的な距離が最小となる時刻を表す部分集合を
\[
  T'' = \{ \tilde{t}_1, \tilde{t}_2, \dots, \tilde{t}_K \} \quad (K \geq 2)
\]
として選択する．
ここで $T''$ は，最多区間内において互いに時間的距離が最も近い2点以上の時刻を選択した集合であり，複数人が同時に帰宅しやすい中心的な時刻を表すものとする．


基準時刻 $\bar{t}$ は，選択された時刻集合 $T''$ の平均として，
\[
  \bar{t} = \frac{1}{K} \sum_{k=1}^{K} \tilde{t}_k
\]
と定義する．
基準時刻を最多区間内で時間的距離が最も近い複数人の平均として定義したのは，「すでに一緒に帰りやすい少人数の集団」を中心に帰宅行動が誘発されると考えたためである．
時間的に近いユーザの群を基準とし，「数人が帰るなら自分も帰ろう」という心理を反映した時刻算出を意図している．


各滞在者の帰宅予測時刻 $t_i$ に対する重み $w_i$ を，
基準時刻 $\bar{t}$ からの時間的距離に基づき次式で定義する．
\[
  w_i =
  \begin{cases}
    1                          & (|t_i - \bar{t}| = 0) \\
    \dfrac{1}{|t_i - \bar{t}|} & (|t_i - \bar{t}| > 0)
  \end{cases}
\]
重みには，基準時刻からの時間的距離の逆数を用いた．
これにより，基準時刻に近いユーザほど帰宅推奨時刻に与える影響が大きくなり，時間的に離れたユーザの影響は緩やかに低減される．
単純な線形減衰や閾値処理ではなく逆数を用いて，遠い時刻の影響を完全に排除せず，集団全体の傾向をなだらかに反映できる設計としている．


以上を用いて，帰宅推奨時刻 $t_{\mathrm{rec}}$ を
重み付き平均として次式により算出する．
\[
  t_{\mathrm{rec}} =
  \frac{\sum_{i=1}^{N} t_i \, w_i}
  {\sum_{i=1}^{N} w_i}
\]
この重み付き平均により，「一緒に帰りやすい中心的な集団」を軸としつつ，周辺のユーザの帰宅時刻も考慮した帰宅推奨時刻の算出を実現している．
単純平均や中央値では，全滞在者の帰宅時刻を一様に扱うが，本研究では同時帰宅の中心となりやすい時刻を基準とし，時間的に近いユーザほど強く反映できるように重み付き平均を用いて帰宅推奨時刻を算出している．


%
% 基準時刻を，最多区間内で互いに近接する複数人の平均として定義したのは，「すでに一緒に帰りやすい少人数の集団」を中心に帰宅行動が誘発されると考えたためである．
% 個々の最頻値や中央値ではなく，時間的に近い利用者群を基準とし，「数人が帰るなら自分も帰ろう」という心理を反映した時刻算出を意図している．

%
% 重みには，基準時刻からの時間的距離の逆数を用いた．
% これにより，基準時刻に近い利用者ほど帰宅推奨時刻に与える影響が大きくなり，時間的に離れた利用者の影響は緩やかに低減される．
% 単純な線形減衰や閾値処理ではなく逆数を用いて，遠い時刻の影響を完全に排除せず，集団全体の傾向をなだらかに反映できる設計としている．

%
% この重み付き平均により，「一緒に帰りやすい中心的な集団」を軸としつつ，周辺の利用者の帰宅時刻も考慮した帰宅推奨時刻の算出を実現している．
